% D.78 candidate. UNREGISTERED, NOTES ONLY, pending full independent audits.
\subsection{A nonzero remaining compact kernel of exact order five}

\begin{statement}{D.78}{Open}{Candidate: the compact kernel at the point minus two}
Take the actual Alexander twist
\[
 s=(\kappa(2),\kappa(3)),\qquad A_s=\mathbb Z_5[[u,v]]C
\]
of D.35. The candidate computes the remaining term of D.77 as
\[
 \mathcal S_s\simeq\mathbb Z/5,\qquad
 \mathcal K_s\simeq\mathbb Z/5,\qquad L_5=\mathcal T_s=0.
                                                               \tag{D.78.1}
\]
The first isomorphism is detected by an explicit class in a
\emph{single finite} Cartier-dual coefficient module, not just by
classes that become locally trivial at different unbounded levels.
Consequently the separated global second cohomology has the exact
topological description
\[
 0\longrightarrow\mathbb Z/5\longrightarrow\mathcal H_s
       \longrightarrow\mathcal L_s\longrightarrow0,     \tag{D.78.2}
\]
with the actual residue space and topology of D.63. No splitting
of (D.78.2) is asserted.

This is a nonempty metabelian section fibre: the rational point
\(t=-2\) has the specified abelian class. Its full metabelian
fibre is a torsor under \(\prod_{\mathbb N}\mathbb Z_5\), by D.74.
The candidate does not determine higher-derived lifting of every
class in that fibre, or the remaining compact kernels for other
abelian pairs.
\end{statement}

\begin{proof}[Candidate proof, pending independent audits]
Write \(\Lambda=\mathbb Z_5[[u,v]]\), \(t_1=1+u,t_2=1+v\),
and \(D=A_s^D\). Use ordinary Pontryagin duality, denoted by a
star, only after taking Cartier duals.
Let \(I_0\) be the actual compact operator image of the untwisted
action \(\rho_0\) in D.35; let \(I_s\) be the actual image of
\(\rho_s\). These groups include their cyclotomic multiplier.
The twist gives a continuous map
\[
 G_{\mathbb Q}\longrightarrow\mathbb Z_5^2\rtimes I_0,
       \qquad g\longmapsto(a(g),b(g),\rho_0(g)),
                                                               \tag{D.78.3}
\]
where \(I_0\) acts on the translation group by \(\chi\).

\paragraph{The global operator image contains every translation.}
The action of a triple in (D.78.3) is
\(t_1^a t_2^b\rho_0(g)\). This representation of its image is
faithful. If that operator is the identity, comparison of its
values on \(1\) and \(t_1\) first gives \(t_1^\chi=t_1\), hence
\(\chi=1\). Restriction to the first and second boundary axes
then gives \(t_1^a=t_2^b=1\), so \(a=b=0\). Its remaining
untwisted operator is the identity as well. Compactness therefore
identifies the actual image with the corresponding closed
subgroup of the semidirect product.

The untwisted action is unramified outside five, by D.39
(equivalently the actual unramified Jacobi-sum image in
\cite[(1.1)--(1.2) and the following paragraph]{DE_NakamuraSakugawaWojtkowiak2017}).
Inertia at two has Kummer vector \((1,0)\), up to a unit
choice of its five-primary tame generator; inertia at three
has vector \((0,1)\). Their closed images therefore contain
both full coordinate translation groups. Put \(H=\mathbb Z_5^2\).
For any element of the projected image \(I_0\), multiply a lift
by a suitable element of \(H\) to remove its two translation
coordinates. Consequently the actual equality is
\[
                         I_s=H\rtimes I_0.              \tag{D.78.4}
\]
This is an equality of compact operator groups; it is not an
assumption of independence of formal arithmetic coefficients.

\paragraph{The unramified image is entirely supplied at five.}
Put \(K=\mathbb Q(\zeta_5)\), let \(\mathfrak p\) be its unique
prime over five, and set \(P=G_{K,\{\mathfrak p\}}^{(5)}\).
Over \(K\), every finite quotient of the actual untwisted
operator image is a five-group, by its positive augmentation
filtration and \(\chi\equiv1\pmod5\). Thus this image factors
through \(P\).

D.1's class-number and unit calculation gives the following
basis of its Kummer Frattini dual:
\[
 H^1(P,\mathbb F_5)
  \simeq\mathcal O_K[1/5]^\times/
                    \mathcal O_K[1/5]^{\times5}
  =\langle\zeta_5,\phi,\pi\rangle_{\mathbb F_5},
 \quad \phi=-\zeta_5^2-\zeta_5^3,\quad \pi=\zeta_5-1.
                                                               \tag{D.78.5}
\]
Here the identification chooses a fifth root of unity, so the
trivial coefficient over \(K\) is identified with \(\mu_5\).
The valuation of \(\pi\) at \(K_{\mathfrak p}\) is one; D.77
proves independence of the two local unit classes
\(\zeta_5,\phi\). Hence localization is injective on (D.78.5).
The closed local decomposition image in the finitely generated
pro-five group \(P\) therefore maps onto its Frattini quotient
and is all of \(P\). The local cyclotomic image also supplies
the full order-four quotient over \(\mathbb Q\).
It follows that
\[
                G_{\mathbb Q_5}\longrightarrow I_0
                         \quad\hbox{is onto}.           \tag{D.78.6}
\]
This argument uses the exact local unit certificate, not a
surjectivity assertion based on rational dimensions.

\paragraph{The exact closed local logarithm line.}
The local completed Kummer group has basis
\(\kappa(5),\kappa(6)\), as in D.75. Set
\[
 \alpha=\frac{\log16}{4\log6},\qquad
 \beta=\frac{\log81}{4\log6}.
                                                               \tag{D.78.7}
\]
They satisfy
\[
 \alpha\equiv2,\qquad \beta\equiv4,\qquad
                    \alpha-\beta\equiv3\pmod5.
                                                               \tag{D.78.8}
\]
Thus all three are units. Locally the two Kummer classes are
\(\alpha\kappa(6),\beta\kappa(6)\). Apply an integral monomial
gauge to make these equalities hold for actual cocycles.
It preserves (D.78.4), the constant Cartier submodule, and the
finite augmentation quotients used below.

Let \(N\) be the kernel of \(G_{\mathbb Q_5}\to I_0\).
The image \(\kappa(6)(N)\) is all of \(\mathbb Z_5\).
If it were zero modulo five, the actual cocycle
\(\kappa(6)\bmod5\) would descend through the surjection
(D.78.6) to \(I_0\). Inflating to \(G_{\mathbb Q,\{5\}}\)
would give a class in
\[
 H^1(G_{\mathbb Q,\{5\}},\mathbb F_5(1))
                       =\mathbb F_5\kappa(5).
 \]
This is impossible locally: valuation distinguishes
\(\kappa(5)\) from the nonzero unit class \(\kappa(6)\).
The latter is nonzero because \(\log6\) has valuation one,
whereas a fifth power of a principal unit has logarithm
divisible by \(25\). Thus the closed additive subgroup
\(\kappa(6)(N)\) contains a unit and is \(\mathbb Z_5\).

It follows that the local operator image has precisely the
primitive translation line
\[
 H_5=\{(\alpha h,\beta h):h\in\mathbb Z_5\}.
 \]
All its other translation coordinates lie on this line as
well, so they can be removed without changing the projected
element of \(I_0\). Therefore, after the specified integral gauge,
\[
                  I_5=\rho_s(G_{\mathbb Q_5})
                            =H_5\rtimes I_0.            \tag{D.78.9}
\]
The common complement in (D.78.4) and (D.78.9) is the actual
untwisted image, with zero translation coordinates.

\paragraph{Continuous translation coinduction and the local restriction.}
Under \(H\), the Cartier dual \(D\) is the group of continuous
functions \(H\to\mathbb Q_5/\mathbb Z_5\), with translation
action. Indeed \(\Lambda\) is the completed group ring of
\(H\); on \(H\), the cyclotomic twist is one. This is a
continuous coinduced module, so
\[
 H^j(H,D)=0\quad(j>0),\qquad
 D^H=\lambda:=\mathbb Q_5/\mathbb Z_5(-1).               \tag{D.78.10}
\]
The character on the last group follows from the augmentation
character \(\chi^2\) on \(A_s/H\), with the single Cartier
twist. There is no ambiguity between this constant functional
and constant functions modulo constants in D.75's different,
augmentation-ideal model.

For clarity, the acyclicity assertion has an elementary
continuous proof. Homogeneous cochains with values in
continuous functions on the compact translation group identify
with continuous functions in one additional group variable.
Inserting the identity in that variable contracts the
positive-degree complex. All functions involved have finite
image locally, and compactness makes the required partitions
uniformly finite, so the contraction is continuous.
The same argument works for \(H_5\): a continuous linear
complement to its primitive line identifies \(H\) with
\(H_5\times\mathbb Z_5\), and the second factor is simply
part of the discrete function coefficient module.

Let \(D_{\rm lin}=D^{H_5}\). It is the Cartier dual of
\[
 B=\Lambda/(t_1^\alpha t_2^\beta-1)
                \simeq\mathbb Z_5[[z]],\qquad t=1+z,
 \quad t_1=t^\beta,\quad t_2=t^{-\alpha}.
                                                               \tag{D.78.11}
\]
These substitutions give an integral isomorphism, since
\(\alpha,\beta\) are units. Continuous Hochschild--Serre,
using the acyclicity just proved, now gives natural isomorphisms
\[
 \begin{split}
 H^1(I_s,D)&=H^1(I_0,\lambda),\\
 H^1(I_5,D)&=H^1(I_0,D_{\rm lin}).
 \end{split}                                               \tag{D.78.12}
\]
Under them restriction is induced by the actual inclusion
\(\lambda\subset D_{\rm lin}\). Neither group on the right
is being inferred from a formal Lie algebra.

\paragraph{Its entire kernel is of exact order five.}
Choose an actual \(\delta\in I_0\) with \(\chi(\delta)=1\)
and primitive \(\kappa_3(\delta)\), supplied locally by D.77.
Its multiplier on \(B\) is \(\mathcal B_\delta(t^\beta,t^{-\alpha})\).
Formula (D.77.9) and (D.78.8) give
\[
 \mathcal B_\delta(t^\beta,t^{-\alpha})-1
      =z^3\,e(z),\qquad e(0)\in\mathbb Z_5^\times .
                                                               \tag{D.78.13}
\]
The leading coefficient is
\[
 e(0)=\frac{\kappa_3(\delta)}{48}\,
                     \beta(-\alpha)(\beta-\alpha).
                                                               \tag{D.78.14}
\]
Thus \(D_{\rm lin}^{I_0}\) is contained in the dual of
\(B/z^3B\). Its three successive characters are
\(\chi^{-1},\chi^{-2},\chi^{-3}\). For an actual element
\(k\) with \(\chi(k)=2\), all three differences from one
are units modulo five. Hence
\[
                         D_{\rm lin}^{I_0}=0.           \tag{D.78.15}
\]

The quotient \(D_{\rm lin}/\lambda\) is the Cartier dual
of the stable ideal \(zB\). By (D.78.13), its invariant
part is contained in
\[
                       V=(zB/z^4B)^D.
 \]
On this module the cyclotomic kernel of \(I_0\) acts
trivially: its exact beta multiplier starts in degree
three, which acts trivially after multiplication by \(z\)
and reduction modulo \(z^4\). The three successive
characters of \(V\) are
\(\chi^{-2},\chi^{-3},\chi^{-4}\).
The three eigenvalues of \(k\) have pairwise unit
differences; its triangular integral matrix is therefore
diagonalized over \(\mathbb Z_5\), by successive elementary
changes of basis. Its invariant subgroup on the torsion
dual has order
\[
       \#\ker(2^{-4}-1:\mathbb Q_5/\mathbb Z_5
                         \to\mathbb Q_5/\mathbb Z_5)=5,
 \]
since \(v_5(2^4-1)=1\); the other two kernels are zero.
The element two topologically generates \(\mathbb Z_5^\times\).
We obtain
\[
                (D_{\rm lin}/\lambda)^{I_0}\simeq\mathbb F_5.
                                                               \tag{D.78.16}
\]
By the cohomology sequence of
\(0\to\lambda\to D_{\rm lin}\to D_{\rm lin}/\lambda\to0\)
and (D.78.15), (D.78.16) identifies the entire kernel of
the restriction map in (D.78.12), not merely a subgroup.

Its generator has a precise arithmetic description.
Let \(c(g)\in\mathbb F_5\) be the reduction of the normalized
integral weight-three cocycle \(48h(g)\) in D.31.
Its restriction to \(\ker\chi\) is the raw \(\kappa_3\),
and D.39 shows that it factors through \(I_0\).
Modulo five its character \(\chi^3\) equals \(\chi^{-1}\),
so it defines a cocycle with values in \(\lambda[5]\).
It is nonzero because its restriction to the local
cyclotomic kernel is primitive.

To identify the connecting homomorphism in (D.78.16),
lift its invariant vector to a functional on \(B/z^4B\)
of order five with nonzero top coefficient. Such a lift
can be made \(k\)-fixed by changing its constant coefficient,
since \(k-1\) is invertible on \(\lambda[5]\).
For every \(\tau\in\ker\chi\), (D.78.14) gives the defect
as a fixed nonzero scalar times
\(\kappa_3(\tau)\) times the constant functional.
Thus its connecting class is a nonzero scalar multiple
of \([c]\). Equality follows from restriction to
\(\ker\chi\), since
\(H^1(\mathbb Z_5^\times,\mathbb F_5(-1))=0\), by the
order-four scalar. All lifts and corrections can be
made in the order-five dual of \(B/z^4B\).
In particular the class \([c]\) has a local primitive at
five in a finite three-jet module.

\paragraph{An upper bound for the actual inverse finite-localization kernel.}
Finite Poitou--Tate duality in D.63 identifies
\[
 \mathcal S_s^*=\varinjlim_r
       \Sha^1(G_{\mathbb Q},W_r^D),
       \qquad W_r=A_s/(5,u,v)^rA_s,                     \tag{D.78.17}
\]
where \(\Sha^1\) is the kernel of restriction to all places
at that \emph{one finite coefficient level}. The direct
limit injects into \(H^1(G_{\mathbb Q},D)\), by exactness
of filtered colimits of the finite-coefficient cohomology
groups and their subgroups. This does not commute a
colimit with an infinite product of local groups.

Every finite class in (D.78.17) descends to the actual
finite operator image. If a cocycle \(z_r\) were nonzero
on an element acting trivially on \(W_r^D\), the finite
affine image \((z_r,\rho_r)\) would contain
\((m,1)\) with \(m\ne0\). Chebotarev at a prime outside
its finite ramification set would realize a conjugate
Frobenius. A local coboundary there evaluates to zero on
that Frobenius, whereas this cocycle does not. This
contradicts its local triviality. Hence \(z_r\) factors
through that finite image, and therefore through \(I_s\).
Inflation from \(I_s\) to the global group is injective
in degree one. Likewise inflation from \(I_5\) to the
local group at five is injective. Thus (D.78.17) injects
into the restriction kernel computed in (D.78.16):
\[
                         \mathcal S_s^*
                                  \lhook\joinrel\longrightarrow\mathbb F_5.
                                                               \tag{D.78.18}
\]
The Chebotarev use is the exact finite-image theorem
\cite[V, Theorem 3.23]{PBC_MilneCFT}.

\paragraph{One bounded coefficient module works at every place.}
It remains to prove that \([c]\) survives in (D.78.17).
This is the uniformity step; pointwise division by \(q-1\)
would not prove it.
Let \(\epsilon_{ij}\) be the coefficient functional on
\(u^iv^j\), with values in \(\mathbb F_5\), and work in
the finite Cartier-dual module
\[
 M=\bigl(A_s/(5,(u,v)^4)A_s\bigr)^D
                  \subset W_4^D.
                                                               \tag{D.78.19}
\]
The constant functional is \(\epsilon_{00}\), and the
global cocycle is \(z(g)=c(g)\epsilon_{00}\).

For \(\bar\chi=\chi(g)\bmod5\ne1\), it has a primitive in
\(\mathbb F_5\epsilon_{00}\), since \(\bar\chi^{-1}-1\)
is a unit. Suppose now that \(\bar\chi=1\).
Through total degree three, D.31/D.39 and the actual
Fox formula give
\[
 \rho_s(g)f=
 t_1^{a(g)}t_2^{b(g)}
 \left(1+\frac{c(g)}{48}uv(u+v)\right)f
       \pmod{5,(u,v)^4}.                                \tag{D.78.20}
\]
Here substitution is the identity to this degree, and
both \(S_\chi\)-factors are one. Indeed their relevant
binomial denominators are units at five and
\(\chi-1\) is divisible by five. The weight-two beta
coefficient \((\chi^2-1)/24\) is zero modulo five;
conversion from logarithmic coordinates therefore
adds no cubic term. This verifies exactly the normalized
coefficient used in (D.78.20).

If \(\bar a\ne0\), then
\[
           (g-1)\epsilon_{10}=-\bar a\,\epsilon_{00},
 \]
and \(-c(g)\bar a^{-1}\epsilon_{10}\) is a primitive.
If \(\bar a=0,\bar b\ne0\), use
\(-c(g)\bar b^{-1}\epsilon_{01}\).
If both are zero, (D.78.20) gives
\[
           (g-1)\epsilon_{21}
                         =-\frac{c(g)}{48}\epsilon_{00},
 \]
so the single vector \(-48\epsilon_{21}\) is a primitive.
These equations use the inverse operator in the dual
action; their minus signs are therefore intentional.
All the primitives are in the one fixed finite module \(M\).

Outside \(\{2,3,5\}\), both the actual module and this
cocycle are unramified. The preceding primitive on
Frobenius consequently trivializes the whole local
cocycle. At two and three the constant primitive works
on the whole decomposition group: its Frobenius
multiplier is respectively two or three modulo five,
inertia fixes the constant functional, and \(c\) is
unramified there. At five the three-jet primitive
constructed after (D.78.16) belongs to \(M\), since
the quotient substitutions (D.78.11) and all integral
monomial gauges preserve the order-five three-jet
module. At the real place odd-primary cohomology is
zero (and the chosen normalized cocycle is zero on
complex conjugation). Thus
\[
                 [c]\in\Sha^1(G_{\mathbb Q},M)
                           \longrightarrow\Sha^1(G_{\mathbb Q},W_4^D).
                                                               \tag{D.78.21}
\]
This is a single global finite-level certificate.

The class remains nonzero in \(H^1(G_{\mathbb Q},D)\).
In (D.78.12) it comes from the nonzero class in
\(H^1(I_0,\lambda[5])\); this class remains nonzero in
\(H^1(I_0,\lambda)\), since \(\lambda^{I_0}=0\) and
the multiplication-by-five coefficient sequence makes
that inclusion injective in degree one. The first
isomorphism of (D.78.12) and global inflation preserve
it. Therefore (D.78.21) gives a nonzero element of
(D.78.17). Together with (D.78.18), it proves
\(\mathcal S_s^*\simeq\mathbb F_5\), hence
\(\mathcal S_s\simeq\mathbb Z/5\).

\paragraph{The exact compact term and actual section fibre.}
Each of \(a,b,a-b\) has global depth zero. By
(D.78.7)--(D.78.8), each also has local-at-five
depth zero. Thus D.77 gives \(L_5=\mathcal T_s=0\)
and \(\mathcal K_s=\mathcal S_s\), proving (D.78.1)
and (D.78.2).

For the actual rational point \(t=-2\), the two
Kummer coordinates are \(\kappa(-2),\kappa(3)\).
The first equals \(\kappa(2)\) in completed pro-five
cohomology, since \(-1\) has zero pro-five Kummer class.
Hence the metabelian fibre is nonempty, and D.35/D.74
give the asserted torsor. There is no assertion that
its nonzero second-cohomology classes obstruct this
particular point section; that section already exists.
\end{proof}
